\section{The primitive and nonlinear-step inputs}\label{sec:inputs}

We use the photon and fermion constructions of
\cite{MaxwellBlockManuscript,FermionBlockManuscript}, the extraction
and reconstruction maps of \cite{NormalizedContraction}, and the
nonlinear maps of \cite{NonlinearStep}.
We recall their hypotheses and the conclusions needed here.
The finite expressions used in the comparison are recalled in \cref{app:formulas}.

For an active step put \(e=e_k\), \(e_{k+1}=\sqrt L e_k\) and
\(0<\mu\le L^{-1}\). Both charges lie below the common ceiling of the
primitive estimates. The incoming and outgoing sides are admissible:
\(\ell/(LM)\) and \(\ell/(L^2M)\) are odd integers at least \(13\).
Assume \(0<\rho<1\), odd \(L\ge4^8\),
\(M=L^m\ge L^2\), \(M>4L\),
\(0<\varepsilon\le\bar\varepsilon=1/200\), and
\(1-\varepsilon<\alpha<1\). The condition \(\ell\ge Mr_k\)
is imposed at every active step.

The constants in the block estimates are
\[
 \begin{gathered}
 c_X=54,\qquad \kappa_0=54\log72,\qquad K_0=\exp(\kappa_0)/36,\qquad
 c_{\rm rb}=c_XK_0,\qquad c_{\rm fl}=61c_{\rm rb},\\
 c_C=3L^{9/2},\qquad b=1,\qquad b_L=b/L,\qquad b_-=1-L^{-1},\\
 C_T=2\exp(3/2)
       [1+b_-^{-1}(54L\exp(6)+12\exp(1)+1)],\\
 c_S=10^5c_X^5K_0^5\exp(6\kappa),\qquad
 \kappa_w=\kappa+\kappa_0+3.
 \end{gathered}
\]
The constants \(C_w,\gamma_w,\alpha_0,C_B\) and \(c_\gamma\) are those
of the primitive families; \(c_\gamma,C_w\ge1\).
Here \(c_\gamma\) bounds the exponential row sums on the photon
coordinate and unit-bond index sets. The determinant constant is
\[
 C_d=16c_XC_w\max\{\tfrac\pi2 b\exp(1)(1+2\exp(1)),
                         2+58\exp(1)+b\exp(1)\}.
\]
For related tori \(C_w\) always means the common constant after
projection, \(64K_0C_w^{(1)}\).

\paragraph{Primitive hypotheses.}
Use the full hypotheses and conclusions of
\cite[Theorems A--B and Proposition 5.1]{MaxwellBlockManuscript}
and \cite[Theorem A and Propositions 4.1--4.3]{FermionBlockManuscript}.
In particular the parameter families are defined on
\(|s_{\rm pc}|\le M^{\alpha_0}\), with common analytic kernel and
short H\"older bounds, exact full-strength values, component
restriction, the row-vanishing properties stated for each family,
and local read bonds.
Write $s=(s',s'')$, where $s'=(s'_M,s'_{LM})$ consists of
the photon reconstruction and square-root parameters, of kinds
$\pi_M,\pi_{LM}$ in \cite[(37)]{MaxwellBlockManuscript}.
The fermion parameters $s''$ are indexed by the partition
$\pi_k$ into $M$-cubes, as in
\cite[\S4.1]{FermionBlockManuscript}.
For a photon site $c$, put
$K_c=\{c':|c'-c|_\infty\le2L\}$; the row of $H_k$ at $(c,g)$
vanishes unless every $M$-cube meeting $K_c$ is active.
Their lattice, gauge, transpose and reality identities are retained.
At zero use the diagonal families and the full chamber covariance.
The local resolvent bound then retains its factor \((1+y)^{-1}\).
The top fermion estimates are taken with polymer decay
\(\kappa_w+\kappa_0\); the projected estimates have decay \(\kappa_w\).

The finite inverse and walk conditions of the fermion construction
and the photon weakening condition
\[
 0<\alpha_0\le\tfrac14,\qquad
 2\alpha_0 n_P\log M\le\gamma_PM,\qquad n_P=402
\]
are assumed. Exact photon intertwining under odd covers is
\cite[(46)]{MaxwellBlockManuscript}. The comparison below proves
intertwining of the exact fermion formulas directly.

\paragraph{Extraction, reconstruction and determinant estimates.}
Use \cite[Theorems A--B and Appendix B]{NormalizedContraction},
including the fine inverses at both signs of the total connection
required by the reconstruction theorem.
Use the weighted integration convention of
\cite[Theorem B and Proposition 5.1]{FermionBlockManuscript}
and the determinant estimates of
\cite[Propositions 6.1--6.2]{FermionBlockManuscript}.
The extraction maps are linear, real on real activities, and satisfy
\[
 |m(\mathscr E)|\le c_{\rm fl}e_k^{1/2}\|\mathscr E\|_k^x,\qquad
 |\varepsilon_{\rm v}(\mathscr E)|\le c_{\rm fl}\|\mathscr E\|_k^x .
\]
The normalized linear term and determinant term satisfy
\[
 \|\mathcal L(R\mathscr E)^{\rm nat}\|_{k+1}^x
      \le c_{\rm fl}L^{-1/4+2\varepsilon}\|\mathscr E\|_k^x,
 \qquad
 \|\mathcal L\delta E_k^{\rm det}\|_{k+1}^x
      \le e_{k+1}^{1/4-2\varepsilon}.
\]
The scaled domains obey
\[
 \mathscr A_L\in L^{-3/4}\widetilde R_k\subset\widetilde R_k^-,
 \qquad \mathscr C_L\in L^{-3/4}\mathfrak X_k.
\]
The complete real activity spaces, all these maps and their norms
commute with restriction to the external slot \(\mathscr C=0\);
this restriction does not increase the norm.

\paragraph{Numerical conditions for the nonlinear step.}
The hypotheses of \cite[Theorem A]{NonlinearStep} include
\[
 \begin{gathered}
 L^{1/8-3\bar\varepsilon/2}\ge2400,\qquad e^\varepsilon\le\tfrac14,\qquad
 160L^{1+3\varepsilon/2}Me^\varepsilon\le1,\\
 9ML^\varepsilon e^{2\varepsilon}\le1,\qquad
 90c_XLMe^\varepsilon\le1,\qquad
 \kappa\ge7\kappa_0+7,\qquad
 L(\kappa-6\kappa_0-4)-1\ge\kappa,
 \end{gathered}
\]
and
\[
 C_w\ge4\exp(2\gamma_w)M^{27\alpha_0},\qquad
 M^{\alpha_0}-1\ge\exp((2L+1)^3\kappa+\log(6c_X)).
\]
For \(v=-\log e\), \(p_{0,k}=v^{p_0}\) and \(p_k=v^p\), impose
\[
 \begin{gathered}
 p_{0,k}\le e^{-\varepsilon/4},\qquad
 12\exp(-p_{0,k}^2/2)\le e^7,\qquad 2e^{5\varepsilon/2}\le1,\qquad
 c_{\rm fl}e^\varepsilon\le1,\qquad 3Me^\varepsilon\le1,\\
 2c_CC_w^2c_\gamma^2e^{\varepsilon/4}\le1,\qquad
 16c_Xc_\gamma C_wC_T^2M^3e^{\varepsilon/2}\le1,\qquad
 96M^3e^{1/4}\le1,\\
 32C_wC_Te^{3\varepsilon/2}\le1,\qquad
 88M^3e^{1/4-6\varepsilon}\le1,\qquad
 112C_w^2c_\gamma e^{2\varepsilon}\le3,\\
 c_Se^{1/4-5\varepsilon}\le1,\qquad c_SL^3e^\varepsilon\le1,\qquad
 8c_{\rm rb}L^3C_dM^3e^{2\varepsilon}\le1,\\
 8C_Bp_k\le e^{-3/4+3\varepsilon},\qquad
 8e^{3/4-7\varepsilon/2}\le1,\qquad L^{3/4}\ge8/3,\qquad
 v\ge9,\qquad 8M(v/4+4)e^\varepsilon\le1,\\
 e^{1/4}+e^\varepsilon<1,\qquad
 4e^{1/4+2\varepsilon}+2e^{3\varepsilon}\le1.
 \end{gathered}
\]
These conditions are imposed at each active scale. The additional
conditions used for the stopping arithmetic and the finite-sequence
contraction are stated with the physical scales.

\paragraph{The imported nonlinear conclusions.}
The displayed finite contour definitions use
\[
 r_1=\tfrac12e^{-3/4+9\varepsilon/2},\qquad
 r_2=e^{-3/4+7\varepsilon},\qquad r_4=e^{-1/4+3\varepsilon}.
\]
For a bond field $Z$ put
$|Z|_* =\max\{\|Z\|_\infty,\|\partial Z\|_\infty,
\|\delta_\alpha\partial Z\|_\infty\}$.
By \cite[Proposition 5.1]{NonlinearStep},
\(2\le r_2,r_4\le r_1\) and
\[
 |Z(s')|_*\le\tfrac12e^{-\varepsilon/2},\qquad
 A+tZ(s')\in\widetilde R_k,\qquad
 A+\mathscr C+tZ(s')\in\widetilde{\mathcal G}_k
       \quad(|t|\le r_1).
\]
The finite activities have the local ordinary, pair and fluctuation
supports and read bonds stated in
\cite[Proposition 5.3]{NonlinearStep}.
Their connected union expansion converges absolutely in the weighted
norm and its exponential is the original finite fluctuation integral,
by \cite[Proposition 6.1]{NonlinearStep}.

For every complex datum in \(\operatorname{Ball}_k^x(4)\), the resulting
correction satisfies
\[
 \|E_k^*(m,\mathscr E)\|_{k+1}^x\le e_k^{1/4-6\varepsilon}.
\]
It is norm analytic on affine complex lines in the interior of the
allowed data ball, continuous on their closed allowed sets, and real
for real data. For two data in \(\operatorname{Ball}_k^x(1)\),
\[
 \|E_k^*(m,\mathscr E)-E_k^*(\widetilde m,\widetilde{\mathscr E})\|_{k+1}^x
 \le e_k^{1/4-6\varepsilon}
 \left(\frac{|m-\widetilde m|}{w_k^{\rm m}}
       +\frac{\|\mathscr E-\widetilde{\mathscr E}\|_k^x}{w_k^{\rm E}}\right).
\]
At \(k=0\) the activity is fixed to zero and the second quotient is
absent. These are \cite[Proposition 8.1]{NonlinearStep}.
The full mass, vacuum, activity and external-scalar recursions reproduce
the exact density and its unnormalized photon-box marginal by
\cite[Theorem A]{NonlinearStep}.
In particular \(0\le\delta\varepsilon_k^0\le e_k^7\).
All these conclusions commute with restriction to \(\mathscr C=0\).


The exact external scalar update, with $\mathfrak P_0=0$, is
\begin{equation}\label{eq:rgnl-external-scalar-map}
 \mathfrak P_{k+1}(Y,\mathscr C')
  =\sum_{\bar X^L=LY}
       [\mathfrak P_k(X,\mathscr C'_L)+E_k^{\mathrm{det}}(X,\mathscr C'_L)].
\end{equation}
The full activity update is
\[
 \mathscr E^+=\mathcal L(R\mathscr E)^{\mathrm{nat}}
                +\mathcal L\delta E_k^{\mathrm{det}}+E_k^*(m,\mathscr E),
 \qquad
 \delta E_k^{\mathrm{det}}(X;A,\mathscr C)
       =E_k^d(X,A+\mathscr C)-E_k^d(X,\mathscr C).
\]
With no external argument the determinant difference is
$E_k^{\mathrm{det}}(X,A)=E_k^d(X,A)-E_k^d(X,0)$.
At the first step the incoming activity vanishes. There are
$192\ell_0^3$ incoming chamber pairs of measure $1/48$ and
$4\ell_1^3$ outgoing pairs. The fluctuation covariance is $48\delta$
on the chamber modes, and the local vertex is
\eqref{eq:rg-step-first-vertex}.
